\subsection*{CU-23: Invitar a un amigo a jugar}
\addcontentsline{toc}{subsection}{CU-23: Invitar a un amigo a jugar}
\label{cu-23}

\begin{fichacu}{CU-23}{Invitar a un amigo a jugar}
  \crudFila{Responsable}{Ángel Gabriel Aguilar Hernández}
  \crudFila{Actor principal}{Jugador que invita}
  \crudFila{Actores interesados}{Amigo al que se invita}
  \crudFila{Actores de sistema}{Servidor de partidas, Base de datos}
  \crudFila{Prototipos}{2a, 10a, 13a}
  \crudFila{Prioridad}{Must (debe estar en la entrega). Categoría (a) Obligatorio: característica 5 del cliente (invitación de amigos), junto con el código de sala que el \mbox{CU-10} permite compartir.}
  \crudFila{Objetivo}{Permitir que el Jugador invite directamente a un amigo conectado a una partida privada, sin que ninguno de los dos tenga que compartir un código.}
  \crudFila{Disparador}{El Jugador selecciona ``invitar a jugar'' junto a un amigo en la \texttt{GUI\_\allowbreak{}Friends}, o ``retar'' en la tarjeta de perfil de un amigo (\mbox{CU-08}).}
  \textbf{Precondiciones} & \cuCelda{\begin{cureglas}
      \item \textbf{\mbox{PRE-01}} El Jugador tiene una \textit{Sesion} abierta y su token es válido.
      \item \textbf{\mbox{PRE-02}} El \textit{Usuario} del Jugador tiene \texttt{tipo\_\allowbreak{}cuenta} en \texttt{REGISTRADA} (\mbox{RN-06}). \textit{(Ver \mbox{FA-10})}
      \item \textbf{\mbox{PRE-03}} El Jugador no participa en ninguna partida en línea sin terminar y no está en la \textit{ColaEmparejamiento} del \mbox{CU-09}.
      \item \textbf{\mbox{PRE-04}} El Servidor de partidas está en ejecución y tiene conexión disponible con la Base de datos.
    \end{cureglas}} \\ \hline
  \textbf{Flujo normal} & \cuCelda{\begin{cuenum}[start=1]
      \item El Jugador selecciona ``invitar a jugar'' junto a un amigo en la \texttt{GUI\_\allowbreak{}Friends}, o ``retar'' en su tarjeta de perfil.
      \item El SJM despliega la hoja rápida de invitación con el \textit{Modo} y el reloj de la última partida privada preseleccionados.
      \item El Jugador confirma los ajustes y selecciona ``Enviar invitación''.
      \item El SJM envía el mensaje \texttt{INVITE\_\allowbreak{}TO\_\allowbreak{}MATCH\_\allowbreak{}REQUEST} con el \texttt{id\_\allowbreak{}usuario} del amigo, los ajustes y el token. \textit{(Ver \mbox{EX-03}, \mbox{EX-04})}
      \item El Servidor de partidas verifica que el token corresponda a una \textit{Sesion} abierta de una cuenta registrada y comprueba la precondición \mbox{PRE-03}. \textit{(Ver \mbox{FA-01}, \mbox{FA-10})}
      \item El Servidor de partidas verifica que exista una \textit{Amistad} entre el Jugador y el destinatario, y que la cuenta del destinatario no esté \texttt{ELIMINADA} (\mbox{RN-02}). \textit{(Ver \mbox{FA-02})}
      \item El Servidor de partidas verifica que el amigo tenga una \textit{Sesion} abierta, que no participe en una partida en línea sin terminar y que no esté en la \textit{ColaEmparejamiento} del \mbox{CU-09} (\mbox{RN-05}). \textit{(Ver \mbox{FA-03}, \mbox{FA-04})}
      \item El Servidor de partidas verifica que no exista ya una \textit{Invitacion} pendiente del Jugador al mismo amigo (\mbox{RN-03}). \textit{(Ver \mbox{FA-05})}
    \end{cuenum}} \\
   & \cuCelda{\begin{cuenum}[start=9]
      \item Si el Jugador no está en una sala abierta, el Servidor de partidas crea la \textit{Sala} con los ajustes, conforme a los pasos 8 a 10 del \mbox{CU-10}; si ya está en una sala propia, usa esa. \textit{(Ver \mbox{EX-01})}
      \item El Servidor de partidas crea la \textit{Invitacion} con la \texttt{id\_\allowbreak{}sala}, el \texttt{id\_\allowbreak{}emisor}, el \texttt{id\_\allowbreak{}destinatario}, la \texttt{fecha\_\allowbreak{}invitacion}, la \texttt{fecha\_\allowbreak{}expiracion} a sesenta segundos y el \texttt{estado} igual a \texttt{PENDIENTE}. \textit{(Ver \mbox{EX-01})}
      \item El Servidor de partidas notifica al amigo con \texttt{MATCH\_\allowbreak{}INVITATION}, con el \texttt{nickname} del Jugador y los ajustes.
      \item El SJM del Jugador despliega la \texttt{GUI\_\allowbreak{}WaitingRoom} de la sala, con la plaza del amigo marcada como ``invitado, esperando respuesta''.
      \item El SJM del amigo muestra un aviso emergente con los ajustes y las opciones ``Aceptar'' y ``Rechazar'', con la cuenta atrás de los sesenta segundos. \textit{(Ver \mbox{FA-06}, \mbox{FA-07})}
      \item El amigo selecciona ``Aceptar'' y su SJM envía \texttt{ACCEPT\_\allowbreak{}INVITATION\_\allowbreak{}REQUEST}.
      \item El Servidor de partidas verifica que la \textit{Invitacion} siga \texttt{PENDIENTE}, no haya expirado y que la sala siga \texttt{ABIERTA} con plaza libre, y que el amigo no haya entrado mientras tanto en una partida o en la \textit{ColaEmparejamiento} del \mbox{CU-09} (\mbox{RN-05}). \textit{(Ver \mbox{FA-08}, \mbox{FA-09})}
      \item El Servidor de partidas marca la \textit{Invitacion} como \texttt{ACEPTADA} y agrega al amigo a la sala conforme a los pasos 9 a 11 del \mbox{CU-11}, sin pedir el código (\mbox{RN-04}). \textit{(Ver \mbox{EX-02})}
    \end{cuenum}} \\
   & \cuCelda{\begin{cuenum}[start=17]
      \item El SJM del amigo despliega la \texttt{GUI\_\allowbreak{}WaitingRoom}, y el flujo continúa en el paso 13 del \mbox{CU-11}.
      \item Fin del caso de uso.
    \end{cuenum}} \\ \hline
  \textbf{Flujos alternos} & \cuCelda{\textbf{\mbox{FA-01}: El Jugador ya está en una partida o en la cola}\par
    \begin{cuenum}[start=1]
      \item El Servidor de partidas responde con \texttt{INVITATION\_\allowbreak{}FAILED} y el motivo \texttt{YA\_\allowbreak{}EN\_\allowbreak{}PARTIDA} o \texttt{YA\_\allowbreak{}EN\_\allowbreak{}COLA}.
      \item El SJM ofrece volver a la partida, que continúa en el \mbox{CU-19}, o cancelar la búsqueda del \mbox{CU-09}.
      \item Fin del caso de uso.
    \end{cuenum}} \\
   & \cuCelda{\textbf{\mbox{FA-02}: El destinatario no es amigo del Jugador o su cuenta no está disponible}\par
    \begin{cuenum}[start=1]
      \item El Servidor de partidas no encuentra la \textit{Amistad} entre ambos, porque se eliminó mientras tanto (\mbox{CU-24}), o encuentra la cuenta del destinatario \texttt{ELIMINADA}.
      \item El Servidor de partidas responde con \texttt{INVITATION\_\allowbreak{}FAILED} y el motivo \texttt{NO\_\allowbreak{}ES\_\allowbreak{}AMIGO} o \texttt{JUGADOR\_\allowbreak{}NO\_\allowbreak{}DISPONIBLE}, sin crear la sala ni la invitación.
      \item El SJM muestra que no es posible invitar a ese jugador y actualiza la lista de amigos.
      \item Fin del caso de uso.
    \end{cuenum}} \\
   & \cuCelda{\textbf{\mbox{FA-03}: El amigo no está conectado}\par
    \begin{cuenum}[start=1]
      \item El Servidor de partidas responde con \texttt{INVITATION\_\allowbreak{}FAILED} y el motivo \texttt{JUGADOR\_\allowbreak{}DESCONECTADO}, sin crear la sala ni la invitación.
      \item El SJM lo indica junto al amigo y recuerda que puede crear una sala y compartirle el código por su cuenta (\mbox{CU-10}).
      \item Fin del caso de uso.
    \end{cuenum}} \\
   & \cuCelda{\textbf{\mbox{FA-04}: El amigo está jugando una partida o en la cola}\par
    \begin{cuenum}[start=1]
      \item El Servidor de partidas responde con \texttt{INVITATION\_\allowbreak{}FAILED} y el motivo \texttt{JUGADOR\_\allowbreak{}EN\_\allowbreak{}PARTIDA} o \texttt{JUGADOR\_\allowbreak{}EN\_\allowbreak{}COLA}, sin crear la sala ni la invitación.
      \item El SJM lo indica junto al amigo.
      \item Fin del caso de uso.
    \end{cuenum}} \\
   & \cuCelda{\textbf{\mbox{FA-05}: Ya hay una invitación pendiente al mismo amigo}\par
    \begin{cuenum}[start=1]
      \item El Servidor de partidas responde con \texttt{INVITATION\_\allowbreak{}ALREADY\_\allowbreak{}PENDING} y los segundos que le quedan.
      \item El SJM lleva al Jugador a la sala donde espera esa invitación.
      \item Fin del caso de uso.
    \end{cuenum}} \\
   & \cuCelda{\textbf{\mbox{FA-06}: El amigo rechaza}\par
    \begin{cuenum}[start=1]
      \item El amigo selecciona ``Rechazar'' y su SJM envía \texttt{DECLINE\_\allowbreak{}INVITATION\_\allowbreak{}REQUEST}.
      \item El Servidor de partidas marca la \textit{Invitacion} como \texttt{RECHAZADA} y notifica al Jugador con \texttt{INVITATION\_\allowbreak{}DECLINED}.
      \item El SJM del Jugador libera la plaza en la sala y lo indica; la sala sigue abierta para invitar a otro amigo o compartir el código.
      \item Fin del caso de uso.
    \end{cuenum}} \\
   & \cuCelda{\textbf{\mbox{FA-07}: La invitación expira sin respuesta}\par
    \begin{cuenum}[start=1]
      \item Transcurren los sesenta segundos sin respuesta.
      \item El Servidor de partidas marca la \textit{Invitacion} como \texttt{EXPIRADA} y notifica a ambos.
      \item El SJM del amigo retira el aviso; el del Jugador libera la plaza.
      \item Fin del caso de uso.
    \end{cuenum}} \\
   & \cuCelda{\textbf{\mbox{FA-08}: La sala dejó de estar disponible antes de aceptar}\par
    \begin{cuenum}[start=1]
      \item El Servidor de partidas encuentra la sala cerrada, llena o en partida.
      \item El Servidor de partidas marca la \textit{Invitacion} como \texttt{EXPIRADA} y responde con \texttt{INVITATION\_\allowbreak{}NO\_\allowbreak{}LONGER\_\allowbreak{}VALID}.
      \item El SJM del amigo muestra que la invitación ya no es válida.
      \item Fin del caso de uso.
    \end{cuenum}} \\
   & \cuCelda{\textbf{\mbox{FA-09}: El amigo entró en la cola o en una partida antes de aceptar}\par
    \begin{cuenum}[start=1]
      \item El Servidor de partidas responde al amigo con \texttt{INVITATION\_\allowbreak{}FAILED} y el motivo \texttt{YA\_\allowbreak{}EN\_\allowbreak{}COLA} o \texttt{YA\_\allowbreak{}EN\_\allowbreak{}PARTIDA}, sin marcar la \textit{Invitacion} (\mbox{RN-05}).
      \item El SJM del amigo ofrece cancelar la búsqueda del \mbox{CU-09} para poder aceptar mientras la invitación no expire; si la cancela, el flujo continúa en el paso 14 del FN.
      \item Si no la cancela, la invitación expira conforme al \mbox{FA-07}.
      \item Fin del caso de uso.
    \end{cuenum}} \\
   & \cuCelda{\textbf{\mbox{FA-10}: El Jugador usa una cuenta de invitado}\par
    \begin{cuenum}[start=1]
      \item El SJM no ofrece ``invitar a jugar'' ni ``retar'' a una cuenta de invitado, que no tiene amigos e indica que solo entra a partidas con el código de una sala (\mbox{RN-06}).
      \item Si aun así llega una solicitud de una cuenta de invitado, el Servidor de partidas la rechaza en el paso 5 con \texttt{INVITATION\_\allowbreak{}FAILED} y el motivo \texttt{CUENTA\_\allowbreak{}INVITADO}, y la registra en la bitácora del servidor.
      \item Fin del caso de uso.
    \end{cuenum}} \\ \hline
  \textbf{Excepciones} & \cuCelda{\textbf{\mbox{EX-01}: Fallo al escribir en la base de datos}\par
    \begin{cuenum}[start=1]
      \item El Servidor de partidas revierte la transacción, de modo que no quede una \textit{Invitacion} hacia una sala que no se creó.
      \item El Servidor de partidas registra la excepción con un identificador de incidencia y responde con \texttt{SERVER\_\allowbreak{}ERROR}.
      \item El SJM muestra la \texttt{GUI\_\allowbreak{}MessageError} con el identificador.
      \item Fin del caso de uso.
    \end{cuenum}} \\
   & \cuCelda{\textbf{\mbox{EX-02}: El amigo acepta y falla la entrada a la sala}\par
    \begin{cuenum}[start=1]
      \item El Servidor de partidas revierte la marca \texttt{ACEPTADA} junto con la entrada fallida.
      \item El SJM del amigo muestra el error y conserva el aviso mientras la invitación no expire.
      \item Fin del caso de uso.
    \end{cuenum}} \\
   & \cuCelda{\textbf{\mbox{EX-03}: Se agota el tiempo de espera de la respuesta}\par
    \begin{cuenum}[start=1]
      \item El SJM no reenvía la invitación y consulta el estado de la sala.
      \item El SJM muestra la sala con la invitación si se creó, o la lista de amigos si no.
    \end{cuenum}} \\
   & \cuCelda{\textbf{\mbox{EX-04}: La conexión se interrumpe}\par
    \begin{cuenum}[start=1]
      \item El SJM muestra la barra de conexión perdida.
      \item Si se pierde la conexión del amigo, la invitación expira conforme al \mbox{FA-07}; si se pierde la del Jugador, la sala se rige por el \mbox{EX-04} del \mbox{CU-10}.
    \end{cuenum}} \\ \hline
  \textbf{Postcondiciones} & \cuCelda{\begin{cureglas}
      \item \textbf{\mbox{POST-01}} Si el caso termina por el FN, la \textit{Invitacion} está \texttt{ACEPTADA} y ambos jugadores son miembros de la misma \textit{Sala} abierta.
      \item \textbf{\mbox{POST-02}} Si el caso termina por el \mbox{FA-06}, el \mbox{FA-07} o el \mbox{FA-08}, la \textit{Invitacion} está \texttt{RECHAZADA} o \texttt{EXPIRADA} y el amigo no entró a la sala.
      \item \textbf{\mbox{POST-03}} No hay más de una \textit{Invitacion} pendiente del mismo emisor al mismo destinatario.
      \item \textbf{\mbox{POST-04}} La sala creada para la invitación se comporta como cualquier sala privada del \mbox{CU-10}, y la partida que salga de ella no modificará el elo ni la \textit{EstadisticaModo} de nadie.
    \end{cureglas}} \\ \hline
  \textbf{Reglas de negocio} & \cuCelda{\begin{cureglas}
      \item \textbf{\mbox{RN-01}} Una invitación a jugar crea o usa una sala privada del \mbox{CU-10}; no es un emparejamiento clasificatorio del \mbox{CU-09}. La partida que salga de ella es amistosa: no cambia el elo ni la \textit{EstadisticaModo} de nadie.
      \item \textbf{\mbox{RN-02}} Solo se puede invitar a un amigo: debe existir una \textit{Amistad} entre ambos. Retar a un jugador que no es amigo queda como requisito deseable (\mbox{DES-23}).
      \item \textbf{\mbox{RN-03}} Una invitación expira a los sesenta segundos, y solo puede haber una pendiente del mismo emisor al mismo destinatario.
      \item \textbf{\mbox{RN-04}} Aceptar una invitación da acceso a la sala sin conocer su código.
      \item \textbf{\mbox{RN-05}} Solo se puede invitar a amigos con sesión abierta que no estén en una partida en línea ni en la \textit{ColaEmparejamiento} del \mbox{CU-09}, y solo pueden aceptar mientras sigan así, porque un jugador no puede estar en una sala y en la cola a la vez.
      \item \textbf{\mbox{RN-06}} Una cuenta de invitado no invita ni recibe invitaciones: no tiene amigos (\mbox{RN-07} del \mbox{CU-21}) y solo entra a partidas con el código de una sala (\mbox{RN-05} del \mbox{CU-03}).
      \item \textbf{\mbox{RN-07}} El sistema no envía invitaciones por correo ni por ningún medio externo. Para jugar con quien no es su amigo o no está conectado, el anfitrión copia el código de su sala y lo comparte por donde quiera (\mbox{CU-10}).
    \end{cureglas}} \\ \hline
  \textbf{Observación} & \cuCelda{\textit{Sobre por qué la invitación usa una sala en lugar de crear la partida directamente.}\par
    Crear la partida en cuanto el amigo acepta sería más rápido, pero dejaría dos formas de empezar una partida privada con reglas distintas. Pasar por la sala hace que la invitación sea solo una forma de entrar sin código: los ajustes, la marca de listo, el chat de espera y el inicio los resuelve el \mbox{CU-10}, y funciona igual para un reto uno contra uno que para completar una partida de cuatro invitando a varios amigos.} \\ \hline
  \textbf{Observación} & \cuCelda{\textit{Sobre las tres formas de entrar a una partida.}\par
    Una partida se alcanza por el emparejamiento clasificatorio (\mbox{CU-09}), por el código de una sala (\mbox{CU-10}, \mbox{CU-11}) o por la invitación de un amigo (este caso). Solo la primera mueve el elo. La versión anterior también invitaba por correo; se quitó porque enviar el código por correo no aporta nada que el jugador no pueda hacer copiándolo, y obligaba a un límite de envíos y a guardar correos ajenos (\mbox{D-29}).} \\ \hline
  \textbf{Datos que toca} & \cuCelda{\begin{cudatos}
      \item \textbf{\textit{Sesion}:} lee la del emisor y comprueba la del destinatario \textit{(pasos 5, 7)}
      \item \textbf{\textit{Usuario}:} lee el \texttt{tipo\_\allowbreak{}cuenta} del emisor y el estado del destinatario \textit{(pasos 5, 6)}
      \item \textbf{\textit{Amistad}:} lee la del par \textit{(paso 6)}
      \item \textbf{\textit{Participacion}, \textit{ColaEmparejamiento}:} lee, para las exclusiones \textit{(pasos 5, 7, 15)}
      \item \textbf{\textit{Invitacion}:} lee las pendientes; crea; escribe \texttt{estado} \textit{(pasos 8, 10, 16, \mbox{FA-06}, \mbox{FA-07}, \mbox{FA-08})}
      \item \textbf{\textit{Sala}, \textit{SalaParticipante}:} crea o usa la sala; agrega al amigo \textit{(pasos 9, 16)}
    \end{cudatos}} \\ \hline
\end{fichacu}
\clearpage
